%============================================================
% Chapter 22: Bit Trees: 64-ary Fast Sets and Packed Structures
% STATUS: PLANNED  (generated by tools/outline.py; replace this file when the chapter is written)
%============================================================
\chapter{Bit Trees: 64-ary Fast Sets and Packed Structures}
\label{ch:bit_trees}

\begin{tcolorbox}[enhanced, colback=white, colframe=secondaryblue, boxrule=1pt, arc=2mm]
\textbf{\large\textcolor{primaryblue}{A tree of 64-bit words answers insert, erase, successor and predecessor in $\bigO{\log_{64} n}$: faster than any balanced tree.}}
\end{tcolorbox}

\chapterfacts{1900}{2700}{4}{8}{30}

\begin{plannedbox}[Planned Contents --- this chapter is not written yet]
\textbf{Part A --- Lessons}
\begin{enumerate}[leftmargin=*, label=\textbf{A\arabic*.}, itemsep=1pt]
    \item The 64-ary bit tree: one bit per child, \texttt{ctz}/\texttt{clz} to descend
    \item Successor and predecessor queries; comparison with \texttt{std::set}
    \item Packed boolean arrays and bit-packed small integers
    \item Van Emde Boas trees as the idea behind it
    \item Lowbit inside Fenwick trees is the tree book
\end{enumerate}
\textbf{Part B --- Patterns} (with candidate problems, checked when the chapter is written)
\begin{enumerate}[leftmargin=*, label=\textbf{P\arabic*.}, itemsep=1pt]
    \item \textbf{Successor and predecessor} --- Predecessor Problem {\small\color{darkgray}(Library Checker)}
    \item \textbf{Packed arrays}
    \item \textbf{Fast sets in sweeps}
    \item \textbf{Bit-tree tricks}
\end{enumerate}
\end{plannedbox}
